%
% 32-komplexj.tex -- Jacobi-Matrix einer komplexen Funktion
%
% (c) 2026 Prof Dr Andreas Müller
%
\subsection{Die komplexe Jacobi-Matrix einer komplexen Funktion}
Die ursprünglichen Ableitungen nach $x$ und $y$ können durch Linearkombination
der beiden partiellen Ableitungsoperatoren $\partial/\partial z$ und
$\partial/\partial\bar{z}$ wiedergewonnen werden:
\begin{align*}
\frac{\partial f}{\partial x}
=
\frac{\partial u}{\partial x}
+
i\frac{\partial v}{\partial x}
&=
\frac{\partial f}{\partial z}
+
\frac{\partial f}{\partial \bar{z}}
&&\text{und}
&
-i
\frac{\partial f}{\partial y}
=
-i\frac{\partial u}{\partial y}
+
\frac{\partial v}{\partial y}
&=
\frac{\partial f}{\partial z}
-
\frac{\partial f}{\partial \bar{z}}
\end{align*}
Die partiellen komplexen Ableitungsoperatoren enthalten also die
gesamte Information, die die partiellen Ableitungen nach den reellen
Variablen $x$ und $y$ wiedergeben.

Auch der Real- und Imaginärteil einer komplexen Zahl
$\Delta z=\Delta x + i\Delta y$
lässt sich gemäss \eqref{buch:holomorph:nichtholomorph:eqn:xyzzbar}
aus $\Delta z$ und $\overline{\Delta z}=\Delta x -i\Delta y$
rekonstruieren lassen.
Die Jacobi-Matrix der Funktion $f$ betrachtet als reelle Funktion
$\mathbb{R}^2\to\mathbb{R}^2$ sollte sich also auch als eine Matrix
ausdrücken lassen, die auf den 2-dimensionalen komplexen
Vektor mit Komponenten $\Delta z$ und $\overline{\Delta z}$ wirkt.

Da wir nicht a priori davon ausgehen möchten, dass die Funktion $f$
holomorph ist, schreiben wir sie als $f(z,\bar{z})$ und untersuchen
später, was sich ergibt, wenn $f$ holomorph ist, also nicht von $\bar{z}$
abhängt.

Die reelle Jacobi-Matrix der Funktion $f(z) = u(x,y) + iv(x,y)$
ermöglicht die Konstruktion der linearen Ersatzfunktion von $f$
betrachtet als Funktion zweier reeller Variablen.
Wir verwenden Vektornotation, um diese wie folgt auszudrücken:
\begin{align*}
\begin{pmatrix}
u(x+\Delta x,y+\Delta y)\\
v(x+\Delta x,y+\Delta y)
\end{pmatrix}
&=
\begin{pmatrix}
u(x,y)\\
v(x,y)
\end{pmatrix}
+
\bgroup
\renewcommand{\arraystretch}{1.9}
\begin{pmatrix}
\displaystyle
\frac{\partial u}{\partial x}(x,y) &
\displaystyle
\frac{\partial u}{\partial y}(x,y) \\
\displaystyle
\frac{\partial v}{\partial x}(x,y) &
\displaystyle
\frac{\partial v}{\partial y}(x,y)
\end{pmatrix}
\egroup
\begin{pmatrix}
\Delta x \\
\Delta y
\end{pmatrix}
+
o(\Delta x,\Delta y)
\end{align*}
In diesem Ausdruck möchten wir alles durch die Funktionen $f(z,\bar{z})$
und ihre konjugiert komplexe Funktion $\bar{f}(z,\bar{z})$ sowie die Variablen
$z$ und $\bar{z}$ und die partiellen Ableitungen nach $z$ und $\bar{z}$
ausdrücken.
Für das Inkrement findet man
\[
\begin{pmatrix}
\Delta x\\
\Delta y
\end{pmatrix}
=
\begin{pmatrix}
\frac12   (\Delta z+\overline{\Delta z}) \\
\frac1{2i}(\Delta z-\overline{\Delta z})
\end{pmatrix}.
\]
Die Ableitungen werden
\begin{align*}
\frac{\partial u}{\partial x}
&=
\frac12
\biggl(
\frac{\partial f}{\partial x}
+
\frac{\partial \bar{f}}{\partial x}
\biggr)
=
\frac12
\biggl(
\frac{\partial f}{\partial z}
+
\frac{\partial f}{\partial \bar{z}}
+
\frac{\partial \bar{f}}{\partial z}
+
\frac{\partial \bar{f}}{\partial \bar{z}}
\biggr),
\\
\frac{\partial u}{\partial y}
&=
\frac12
\biggl(
\frac{\partial f}{\partial y}
+
\frac{\partial \bar{f}}{\partial y}
\biggr)
=
\frac12
\biggl(
i
\biggl(
\frac{\partial f}{\partial z}
-
\frac{\partial f}{\partial\bar{z}}
\biggr)
+
i
\biggl(
\frac{\partial\bar{f}}{\partial z}
-
\frac{\partial\bar{f}}{\partial\bar{z}}
\biggr)
\biggr),
\\
\frac{\partial v}{\partial x}
&=
\frac1{2i}
\biggl(
\frac{\partial f}{\partial x}
-
\frac{\partial\bar{f}}{\partial x}
\biggr)
=
\frac1{2i}
\biggl(
\frac{\partial f}{\partial z}
+
\frac{\partial f}{\partial \bar{z}}
-
\frac{\partial \bar{f}}{\partial z}
-
\frac{\partial \bar{f}}{\partial \bar{z}}
\biggr),
\\
\frac{\partial v}{\partial y}
&=
\frac1{2i}
\biggl(
\frac{\partial f}{\partial y}
-
\frac{\partial\bar{f}}{\partial y}
\biggr)
=
\frac1{2i}
\biggl(
i\biggl(
\frac{\partial f}{\partial z}
-
\frac{\partial f}{\partial\bar{z}}
\biggr)
-i
\biggl(
\frac{\partial\bar{f}}{\partial z}
-
\frac{\partial\bar{f}}{\partial\bar{z}}
\biggr)
\biggr).
\end{align*}
Im Interesse etwas kompakterer Formeln schreiben wir die partiellen
Ableitungen nach $z$ und $\bar{z}$ in der in der Theorie der partiellen
Differentialgleichungen üblichen Form durch einen Index in der Form
\begin{align*}
\frac{\partial f}{\partial z}
&=
f_z
&\text{und}&
&
\frac{\partial f}{\partial \bar{z}}
&=
f_{\bar{z}}.
\end{align*}
Damit können wir jetzt die lineare Ersatzfunktion als
\begin{align*}
u(x+\Delta x,y+\Delta y)
&\approx
\frac{\partial u}{\partial x}\Delta x
+
\frac{\partial u}{\partial y}\Delta y
\\
&=
\frac12
(
f_z
+
f_{\bar{z}}
+
\bar{f}_z
+
\bar{f}_{\bar{z}}
)
\frac12
(\Delta z + \overline{\Delta z})
+
\frac{i}2
(
f_z
-
f_{\bar{z}}
+
\bar{f}_z
-
\bar{f}_{\bar{z}}
)
\frac1{2i}
(\Delta z - \overline{\Delta z}),
\\
&=
\frac12
(
f_z + \bar{f}_z
)
\Delta z
+
\frac12
(
f_{\bar{z}} + \bar{f}_{\bar{z}}
)
\overline{\Delta z}
\\
v(x+\Delta x,y+\Delta y)
&\approx
\frac{\partial v}{\partial x}\Delta x
+
\frac{\partial v}{\partial y}\Delta y
\\
&=
\frac1{2i}
(
f_z + f_{\bar{z}} - \bar{f}_z - \bar{f}_{\bar{z}}
)
\frac12
(\Delta z + \overline{\Delta z})
+
\frac12
(
f_z - f_{\bar{z}} - \bar{f}_z + \bar{f}_{\bar{z}}
)
\frac1{2i}
(\Delta z - \overline{\Delta z})
\\
&=
\frac1{2i}
(
f_z - \bar{f}_z
)
\Delta z
+
\frac1{2i}
(
f_{\bar{z}}
-
\bar{f}_{\bar{z}}
)
\overline{\Delta z}
\intertext{schreiben, wobei dass Approximationszeichen $\approx$
immer eine Approximation bis auf einen $o$-Term bedeutet.
Die Funktionen $u$ und $v$ sind Real- und Imaginärteil der
Funktion $f$, daher kann man auch}
f(z+\Delta z, \bar{z}+\bar{\Delta z})
&=
u(x+\Delta x, y + \Delta y)
+
iv(x+\Delta x, y + \Delta y)
\\
&=
f_z\Delta z + f_{\bar{z}}\overline{\Delta z}
+
o(\Delta z, \overline{\Delta z})
\\
\bar{f}(z+\Delta z, \bar{z}+\overline{\Delta z})
&=
u(x+\Delta x, y + \Delta y)
-
iv(x+\Delta x, y + \Delta y)
\\
&=
\bar{f}_z \Delta z
+
\bar{f}_{\bar{z}}\Delta z
+
o(\Delta z, \overline{\Delta z})
\end{align*}
schreiben.
In Matrixform ist dies
\begin{align*}
\begin{pmatrix}
     f (z + \Delta z, \bar{z}+\overline{\Delta z})\\
\bar{f}(z + \Delta z, \bar{z}+\overline{\Delta z})
\end{pmatrix}
&=
\bgroup
\renewcommand{\arraystretch}{2.2}
\begin{pmatrix}
\displaystyle\frac{\partial f}{\partial z}(z,\bar{z}) &
\displaystyle\frac{\partial f}{\partial \bar{z}}(z,\bar{z}) \\
\displaystyle\frac{\partial \bar{f}}{\partial z}(z,\bar{z}) &
\displaystyle\frac{\partial \bar{f}}{\partial \bar{z}}(z,\bar{z})
\end{pmatrix}
\begin{pmatrix}
\Delta z\\
\overline{\Delta z}
\end{pmatrix}.
\egroup
\intertext{Die Matrix auf der rechten Seite kann man als eine
komplexe Form der Jacobimatrix sehen.
Wir können die rechte Seite auch als}
\frac{\partial(f,\bar{f})}{\partial(z,\bar{z})}
\begin{pmatrix}
\Delta z \\
\overline{\Delta z}
\end{pmatrix}
\end{align*}
schreiben.

Für eine holomorphe Funktion wird die komplexe Jacobi-Matrix zu
einer Diagonalmatrix.
Man beachte auch, dass
wegen \eqref{buch:holomorph:nichtholomorph:eqn:dkonj}
die Elemente der unteren Zeile der komplexen Jacobi-Matrix
durch Konjugation und Vertauschen aus den Elementen der oberen
Zeile hervorgehen.



